% =============================================================================
%  overview.tex  —  Template Overview
%
%  This section appears as the first section of the demo document.
%  It explains the project structure and usage conventions to new users.
%  Feel free to delete this file once you are familiar with the template.
% =============================================================================

\section{Template Overview}

This document is a demo for the \textbf{XMUM Assignment Template} —
a general-purpose \LaTeX{} template for English assignments at
Xiamen University Malaysia (XMUM).
It has been maintained and used across courses since May 2025,
covering reports of all sizes from short weekly write-ups to a
500\,+ page MPU4.1 Community Service Final Report.

Features provided by \texttt{xmum\_asg.cls}:

\begin{itemize}
    \item XMUM-style header and ``Page \textit{X} of \textit{Y}'' footer
    \item APA citation style (\verb|\citet{}| / \verb|\citep{}|) via \texttt{apacite}
    \item Syntax-highlighted code blocks
    \item Smart cross-references with \verb|\cref{}| / \verb|\Cref{}| (\texttt{cleveref})
    \item Figure and table numbering per section (Figure~1.1, Table~2.3, \ldots)
    \item Cover page / Rubric / Turnitin report attachment via \verb|\includepdf|
    \item \texttt{[turnitin]} class option — suppress headers/footers for Turnitin submission
\end{itemize}

\subsection{Recommended Project Structure}

Organise your assignment folder as shown below.
This layout separates source content, images, and university documents
into clearly defined locations.

\begin{verbatim}
your-assignment/
├── main.tex                 <- Entry point; edit metadata, then \input sections
├── xum_review.cls           <- Template class; do not edit unless necessary
├── ref.bib                  <- BibTeX references (APA format)
│
├── sections/                <- One .tex file per section (strongly recommended)
│   ├── 01_introduction.tex
│   ├── 02_methodology.tex
│   └── ...
│
├── figure/                  <- All images go here
│   └── my_chart.png
│
└── docs/                    <- Official documents from the university
    ├── cover_page.pdf       <- Exported from cover_page.docx
    ├── sample_rubric.pdf    <- Marking rubric from the course portal
    └── turnitin_report.pdf  <- Downloaded after Turnitin submission
\end{verbatim}

\subsection{Modular Section Management}

\textbf{Strongly recommended:} split your content into one \texttt{.tex} file
per section and \verb|\input| them from \texttt{main.tex}:

\begin{verbatim}
\input{sections/01_introduction.tex}
\input{sections/02_methodology.tex}
\input{sections/03_results.tex}
\end{verbatim}

Do \emph{not} write all your content directly inside \texttt{main.tex}.

\textbf{Why?} This approach makes AI-assisted writing significantly more
effective.
AI tools such as ChatGPT, Claude, and Gemini operate within a limited
context window.
When your entire document is one long file, the AI can only see a portion of
it at a time, often missing critical context from earlier or later sections.
By feeding one section file at a time, the AI reads your complete section in
full, produces higher-quality suggestions, and edits more accurately without
wasting its context budget on unrelated content.

\subsection{Compilation}

This template requires \textbf{XeLaTeX} (not pdfLaTeX).
The number of passes depends on whether your document contains citations.

\textbf{With citations} (i.e.\ you use \verb|\bibliography{ref}|):

\begin{verbatim}
xelatex main   (1st pass — generates .aux)
bibtex main    (resolves citations from ref.bib)
xelatex main   (2nd pass — inserts bibliography)
xelatex main   (3rd pass — resolves all cross-references)
\end{verbatim}

\textbf{Without citations} (no \verb|\bibliography{}| in your document):

\begin{verbatim}
xelatex main   (1st pass)
xelatex main   (2nd pass — resolves cross-references and ToC)
\end{verbatim}

In \textbf{VS Code} with the LaTeX Workshop extension, set the recipe to
\texttt{latexmk} — it detects citations automatically and runs the correct
number of passes without manual intervention.
